% ECONOMICS_PROBLEM: {"id": "J71-531", "title": "Mechanisms of discrimination", "jel_code": "J71", "source_id": "OP-0026"}
% !TeX program = xelatex
\documentclass[11pt,a4paper]{article}
\usepackage[margin=23mm]{geometry}
\usepackage{fontspec}
\setmainfont{Latin Modern Roman}
\usepackage{amsmath,amssymb,mathtools}
\usepackage{xcolor,enumitem,booktabs,longtable,array,xurl,hyperref}
\definecolor{muted}{HTML}{53636C}
\hypersetup{colorlinks=true,urlcolor=blue,linkcolor=blue}
\setlength{\parindent}{0pt}
\setlength{\parskip}{5pt}
\newcommand{\R}{\mathbb R}
\newcommand{\E}{\mathbb E}
\newcommand{\N}{\mathbb N}
\newcommand{\ind}{\mathbf 1}
\newcommand{\Var}{\operatorname{Var}}
\newcommand{\Cov}{\operatorname{Cov}}
\newcommand{\supp}{\operatorname{supp}}
\DeclareMathOperator*{\argmin}{arg\,min}
\DeclareMathOperator*{\argmax}{arg\,max}
\newcommand{\entrymeta}[3]{{\sffamily\small\color{muted}\textbf{\texttt{#1}}\quad|\quad #2\par #3\par}}
\newcommand{\Problemref}[1]{\texttt{#1}}
\newcommand{\JELref}[2]{\texttt{#2} (JEL \texttt{#1})}
\begin{document}
\section*{Mechanisms of discrimination}
\textbf{Problem ID:} \texttt{J71-531}\quad \textbf{JEL:} \texttt{J71}\par
% BEGIN ECONOMICS STATEMENT
\label{problem:OP-0026}
\label{atlas:prob:L6}
\noindent\textbf{Original atlas ID:} L6 (entry 26). \textbf{Classification:} Editorial mechanism identification and career-policy problem.

\noindent\textbf{Original atlas question:} How much disparity comes from tastes, statistical inference, networks, evaluation rules or accumulated disadvantage?

\subsection*{Definitions and assumptions}
Fix a hiring stage, employer population, and distribution $G$ of standardized application characteristics $X$. Let $A\in\{0,1\}$ be a randomized identity cue, such as a name convention that changes perceived group membership, and let $B$ specify an independently randomized information or evaluation regime. Denote the potential callback by $C_j(a,b,x)\in\{0,1\}$. Assume randomized assignment, consistency, positive assignment probabilities, and the announced interference restriction. A cue may convey multiple perceived attributes; interpreting it solely as group identity requires additional evidence. For careers, define discounted earnings $Y_i(\pi)$ under a specified employer policy $\pi$, including nonhire and nonemployment outcomes.
\subsection*{Formal research question}
Identify the standardized callback contrast
\[
 \Delta(b)=\int \mathbb E_j[C_j(1,b,x)-C_j(0,b,x)]\,dG(x)
\]
and its change across evaluation regimes. Then, for an explicit class $\mathcal M$ of employer-belief, preference, network, and household-history models, define
\begin{align*}
 \vartheta(M)&=(\theta_{\mathrm{taste}},\theta_{\mathrm{belief}},\theta_{\mathrm{network}}),\\
 \chi(M)&=\mathbb E_M[Y_i(\pi_1)-Y_i(\pi_0)],\\
 \mathcal T(P)&=\{(\vartheta(M),\chi(M)):M\in\mathcal M(P)\}.
\end{align*}
Here $P$ is the observed joint law of experimental and career measurements, $\mathcal M(P)$ consists of models in $\mathcal M$ matching that law, and each $\theta$ is an announced structural parameter, not a share inferred merely from an outcome gap. Characterize $\mathcal T(P)$. Determine which factorial experiments, belief elicitations, quality signals, or follow-up designs distinguish models with the same callback contrasts but different policy responses.
\subsection*{Interpretation and scope}
Randomized cues can establish differential treatment in the tested setting. They do not by themselves partition it into taste-based and statistical mechanisms: unobserved employer beliefs and cue interpretations can generate observational equivalence. Revealing validated quality can help only under assumptions about employers' comprehension and use of the signal. Blind evaluation can change both information and attention. Conditioning subsequent wage analysis on successful hiring creates treatment-dependent selection; unconditional baseline-cohort outcomes or an explicitly justified selection model are required. The total earnings disparity between actual groups also includes unequal opportunity histories and application composition, which the standardized résumé experiment intentionally holds fixed. Interventions on networks or evaluation rules must be defined separately to measure those pathways.
\subsection*{Source audit and status}
[S1] studies orchestral screening, while [S2] and [S3] provide correspondence-experiment evidence. Their findings are existing evidence, not unresolved questions. The mechanism-identification and career-policy program is an editorial synthesis with an explicitly restricted model class. The sources do not prove a single cause of discrimination or certify universal unresolved status as of 2026.

\subsection*{References}
\begin{enumerate}[label=\textnormal{[S\arabic*]},leftmargin=*]
\item Claudia Goldin and Cecilia Rouse (2000). \emph{Orchestrating Impartiality: The Impact of Blind Auditions on Female Musicians}. American Economic Review 90(4), 715--741. \url{https://doi.org/10.1257/aer.90.4.715}. \textit{Locator:} Primary publisher, working-paper, or author record: bibliographic information and abstract; full-text theorem verification is not claimed. \textit{Role:} Evidence on screened identity in auditions; a specific evaluation setting.
\item Marianne Bertrand and Sendhil Mullainathan (2004). \emph{Are Emily and Greg More Employable Than Lakisha and Jamal? A Field Experiment on Labor Market Discrimination}. American Economic Review 94(4), 991--1013. \url{https://doi.org/10.1257/0002828042002561}. \textit{Locator:} Primary publisher, working-paper, or author record: bibliographic information and abstract; full-text theorem verification is not claimed. \textit{Role:} Randomized résumé-name callback evidence; supports unequal treatment, not a unique mechanism.
\item Patrick Kline, Evan K. Rose, and Christopher R. Walters (2022). \emph{Systemic Discrimination Among Large U.S. Employers}. The Quarterly Journal of Economics 137(4), 1963--2036. \url{https://academic.oup.com/qje/article-abstract/137/4/1963/6605934}. \textit{Locator:} Primary publisher, working-paper, or author record: bibliographic information and abstract; full-text theorem verification is not claimed. \textit{Role:} Employer-level correspondence experiment and heterogeneity in contact gaps.
\end{enumerate}

\subsection*{Common conventions: Source mathematical conventions}

Unless an entry states otherwise, agent and firm sets are finite. State and action spaces are measurable, strategies and outcome maps are measurable, probability laws are specified, and expectations exist for the targets under discussion. Infinite-horizon discount factors lie in $(0,1)$ and discounted objectives are finite. Norms, tolerances and complexity input sizes must be specified for computational comparisons. Constants or primitives that appear locally are inputs, not empirical facts asserted by this catalogue.

\textbf{Identification.} A statistical model $m\in\mathcal M$ implies an observable law $P_m$ and target $\theta(m)$. At law $P$, its identified set is
\[
\Theta(P;\mathcal M)=\{\theta(m):m\in\mathcal M,\ P_m=P\}.
\]
Point identification holds when this set is a singleton, or equivalently when $P_{m_1}=P_{m_2}$ implies $\theta(m_1)=\theta(m_2)$ for all compatible models. Sharp bounds mean bounds attained, or approached when the boundary is not attained, by compatible models. Writing this set is a definition, not a solution: the substantive task is to establish informative restrictions and derive or estimate the set. An unrestricted-confounding model may give only trivial bounds.

\textbf{Inference.} Identification, estimability and valid uncertainty quantification are separate questions. Fix $\alpha\in(0,1)$. A confidence procedure $C_n$ over a declared law class $\mathcal P$ has asymptotic uniform coverage at level $1-\alpha$ when
\[
\liminf_{n\to\infty}\inf_{P\in\mathcal P}\Pr_P\{\theta(P)\in C_n\}\geq1-\alpha.
\]
For set-identified targets the coverage event must be defined separately, for example coverage of the full identified set. If an entry uses identification-indexed models instead of uniquely indexed laws, the same coverage convention is applied over those models. Pointwise limits do not establish uniform validity, and asymptotic validity does not guarantee finite-sample precision. Sample sizes, dependence structures and weak-identification sequences are part of each application.

\textbf{Counterfactuals.} $Y_i(a)$ is a potential outcome under a specified intervention, including its timing, implementation and versions. It is not $Y_i$ conditional on voluntarily choosing $a$. A full intervention vector permits interference; shorthand scalar treatment notation does not silently impose no interference. Assignment, selection, attrition, anticipation and measurement must be specified. A claim about an instrument's local effect is not automatically a population policy effect.

\textbf{Economic equilibrium.} A counterfactual model includes best responses, constraints, feasibility, market clearing, state transitions and expectations consistent with the stated information. When several equilibria exist, either a declared selector is an input or the target is set-valued. Comparisons across policy regimes must use the stated selection convention. Reduced-form relationships are not presumed invariant after an equilibrium-changing intervention.

\textbf{Optimization and welfare.} Optimization is over the stated feasible set. If attainment is not established, $\sup$ or $\inf$ and $\varepsilon$-optimal choices replace an asserted maximizer or minimizer. For example, with value $V=\sup_{a\in\mathcal A}W(a)<\infty$, an $\varepsilon$-optimal set is $\{a\in\mathcal A:W(a)\geq V-\varepsilon\}$ for $\varepsilon>0$. Benefits, costs, risk and health or environmental outcomes can be aggregated only using declared units and conversion weights. Interpersonal utility weights, the treatment of fiscal transfers and foreign profits, discounting, and the behavioral welfare criterion are explicit inputs. A comparative-static derivative additionally requires existence and appropriate differentiability of the selected counterfactual.

\textbf{Sources.} Local labels [S1], [S2], and so forth restart in every question. Locators identify the relevant abstract, model, section or result at the level actually checked; a bibliographic or abstract-level verification is not represented as inspection of an entire proof. Working papers are distinguished from later publications and changing versions. Source summaries are paraphrased. All URL checks and editorial formulations refer to the audit date stated above.
% END ECONOMICS STATEMENT
\end{document}
